% Generado por scripts/analizar_encuesta.py — no editar a mano.

\begin{table}[t]
  \centering\small
  \caption{Frecuencias de las preguntas de opción única. P1 a P3 sobre las 152 respuestas; P4 en adelante, sobre las 116 del segmento. Fuente: elaboración propia.}
  \label{tab:enc-unica}
  \begin{tabular}{l >{\raggedright\arraybackslash}p{8.4cm} rr}
    \hline
    \textbf{Pregunta} & \textbf{Respuesta} & \textbf{Casos} & \textbf{\%} \\
    \hline
    P1 (n = 152) & Menos de 18 & 3 & 2,0 \\
     & 18-25 & 89 & 58,6 \\
     & 26-33 & 42 & 27,6 \\
     & 34-40 & 10 & 6,6 \\
     & Más de 40 & 8 & 5,3 \\
    \hline
    P2 (n = 152) & Varias veces al día & 52 & 34,2 \\
     & Una vez al día & 42 & 27,6 \\
     & Algunas veces por semana & 41 & 27,0 \\
     & Rara vez & 13 & 8,6 \\
     & No la uso & 4 & 2,6 \\
    \hline
    P3 (n = 152) & Sí, frecuentemente & 60 & 39,5 \\
     & A veces & 70 & 46,1 \\
     & No & 22 & 14,5 \\
    \hline
    P4 (n = 116) & Muy frecuentemente & 50 & 43,1 \\
     & Frecuentemente & 37 & 31,9 \\
     & A veces & 16 & 13,8 \\
     & Rara vez & 7 & 6,0 \\
     & Nunca & 6 & 5,2 \\
    \hline
    P5 (n = 116) & Sí, lo compartí & 41 & 35,3 \\
     & Estuve a punto pero no lo hice & 31 & 26,7 \\
     & No & 30 & 25,9 \\
     & No sé & 14 & 12,1 \\
    \hline
    P9 (n = 116) & Menos de 1 minuto & 41 & 35,3 \\
     & 1-3 minutos & 47 & 40,5 \\
     & Más de 3 minutos & 20 & 17,2 \\
     & No verifico & 8 & 6,9 \\
    \hline
  \end{tabular}
\end{table}

\begin{table}[t]
  \centering\small
  \caption{Frecuencias de las preguntas de opción múltiple. Los porcentajes se calculan sobre las 116 respuestas del segmento y suman más de cien. Fuente: elaboración propia.}
  \label{tab:enc-multiple}
  \begin{tabular}{l >{\raggedright\arraybackslash}p{8.4cm} rr}
    \hline
    \textbf{Pregunta} & \textbf{Respuesta} & \textbf{Casos} & \textbf{\%} \\
    \hline
    P8 (n = 116) & Busco en Google & 95 & 81,9 \\
     & Miro la cuenta que lo publicó & 81 & 69,8 \\
     & Comparo con medios que conozco & 55 & 47,4 \\
     & Le pregunto a alguien & 46 & 39,7 \\
     & Nada, sigo de largo & 7 & 6,0 \\
     & Leo los comentarios o las notas de la comunidad si hay & 1 & 0,9 \\
     & Lo busco en X con otras palabras & 1 & 0,9 \\
    \hline
    P12 (n = 116) & Que no tenga sesgo político & 96 & 82,8 \\
     & Que muestre las fuentes/evidencia & 77 & 66,4 \\
     & Que explique por qué llega a ese resultado & 57 & 49,1 \\
     & Que sea transparente sobre su margen de error & 43 & 37,1 \\
     & Que la respalde una institución reconocida & 35 & 30,2 \\
     & Que sea open source & 4 & 3,4 \\
    \hline
    P13 (n = 116) & Que se equivoque (falsos positivos) & 75 & 64,7 \\
     & Que censure opiniones o sátira & 65 & 56,0 \\
     & Sesgo político & 57 & 49,1 \\
     & Privacidad de mis datos & 39 & 33,6 \\
     & Ninguna & 3 & 2,6 \\
     & Que se vuelva paga & 3 & 2,6 \\
     & Que ande lento & 1 & 0,9 \\
    \hline
  \end{tabular}
\end{table}

\begin{table}[t]
  \centering\small
  \caption{Distribución de las escalas de 1 a 5, en porcentaje sobre las 116 respuestas del segmento. Fuente: elaboración propia.}
  \label{tab:enc-likert}
  \begin{tabular}{lrrrrrr}
    \hline
    \textbf{Pregunta} & \textbf{1} & \textbf{2} & \textbf{3} & \textbf{4} & \textbf{5} & \textbf{Media} \\
    \hline
    P6. Preocupación & 0,0 & 3,4 & 14,7 & 37,9 & 44,0 & 4,22 \\
    P7. Capacidad autopercibida & 2,6 & 14,7 & 27,6 & 33,6 & 21,6 & 3,57 \\
    P10. Intención de uso & 2,6 & 6,0 & 13,8 & 35,3 & 42,2 & 4,09 \\
    P11. Confianza en el puntaje & 6,0 & 12,1 & 24,1 & 33,6 & 24,1 & 3,58 \\
    \hline
  \end{tabular}
\end{table}

\begin{table}[t]
  \centering\small
  \caption{Temas de la pregunta abierta (P14), agrupados por respuestas de texto idéntico o de igual pedido. Fuente: elaboración propia.}
  \label{tab:enc-temas}
  \begin{tabular}{lrr}
    \hline
    \textbf{Tema} & \textbf{Total (n = 38)} & \textbf{Segmento (n = 29)} \\
    \hline
    Detectar imágenes generadas con IA & 4 & 3 \\
    Avisar si un contenido viejo recircula como nuevo & 3 & 3 \\
    Mostrar las fuentes o la evidencia & 3 & 2 \\
    No recargar la pantalla & 3 & 1 \\
    Contrastar más de una fuente o postura & 2 & 2 \\
    Distinguir opinión o sátira de un dato falso & 2 & 2 \\
    Ahorrar el tiempo de buscar en Google & 2 & 2 \\
    Avisar de títulos \textit{clickbait} & 2 & 2 \\
    Efectividad & 2 & 2 \\
    Facilidad de uso & 2 & 1 \\
    Otras (una mención cada una) & 7 & 3 \\
    Sin un pedido concreto & 6 & 6 \\
    \hline
  \end{tabular}
\end{table}
